\documentclass[11pt,a4paper]{article}
\usepackage[utf8]{inputenc}
\usepackage[T1]{fontenc}
\usepackage[margin=2.5cm]{geometry}
\usepackage{booktabs}
\usepackage{longtable}
\usepackage{enumitem}
\usepackage{xcolor}
\usepackage{hyperref}
\usepackage{titlesec}
\usepackage{fancyhdr}
\usepackage{tabularx}
\usepackage{amssymb}

\hypersetup{
  colorlinks=true,
  linkcolor=blue!70!black,
  urlcolor=blue!70!black,
}

\pagestyle{fancy}
\fancyhf{}
\fancyhead[L]{\textbf{AETHER-OS}}
\fancyhead[R]{Feature Roadmap \& Architecture}
\fancyfoot[C]{\thepage}

\titleformat{\section}{\Large\bfseries\color{black!85}}{\thesection}{1em}{}
\titleformat{\subsection}{\large\bfseries\color{black!75}}{\thesubsection}{1em}{}
\titleformat{\subsubsection}{\normalsize\bfseries\color{black!65}}{\thesubsubsection}{1em}{}

\title{
  \vspace{2cm}
  {\Huge \textbf{AETHER-OS}}\\[0.5cm]
  {\Large Feature Roadmap: From App to Personal OS}\\[0.3cm]
  {\large \textit{What is needed. Why. What purpose. What you get.}}\\[1cm]
}
\author{Ekin --- Medieninformatik, Hochschule der Medien Stuttgart}
\date{July 2026}

\begin{document}
\maketitle
\thispagestyle{empty}
\newpage

\tableofcontents
\newpage

% =========================================================
\section{Vision Statement}
% =========================================================

AETHER-OS is not a note app with AI bolted on. It is a \textbf{local-first personal operating system} --- the home base for your knowledge, your projects, your tools, and your AI. Everything you need to think, build, and work, in one place, on your machine.

The goal is not to build ``another productivity app.'' The goal is to build the \textbf{first app you open in the morning and the last you close at night} --- a system that replaces the fragmented stack of Activity Monitor, Notes, Todo apps, Calendar, Spotlight, Clipboard managers, Terminal, ChatGPT, and Obsidian with one unified, private, local-first workspace.

% =========================================================
\section{Current State (v0.1)}
% =========================================================

\begin{table}[h]
\centering
\begin{tabularx}{\textwidth}{Xc}
\toprule
\textbf{Feature} & \textbf{Status} \\
\midrule
AI Agent Chat (Ollama, streaming, note context) & \checkmark \\
Semantic Search (vector embeddings) & \checkmark \\
Knowledge Graph (wikilink visualization) & \checkmark \\
Dashboard (vault stats, note preview) & \checkmark \\
AETHER Notes (AI-generated note persistence) & \checkmark \\
Project Dashboard (git status, scan, open in editor) & \checkmark \\
Persistent AI Memory (facts, conversations) & \checkmark \\
Command Bar (Cmd+K palette) & \checkmark \\
Configurable Default Editor & \checkmark \\
Resizable Panels & \checkmark \\
\bottomrule
\end{tabularx}
\end{table}

% =========================================================
\section{Phase 1: The Core OS Layer}
% =========================================================

\subsection{1.1 --- Built-in Terminal}

\textbf{What:} A real PTY-backed terminal inside AETHER-OS. Multiple tabs, split panes, full ANSI support.

\textbf{Why:} You shouldn't leave AETHER-OS to run \texttt{git}, \texttt{npm}, \texttt{cargo}, or \texttt{ollama}. The terminal is the most fundamental developer tool --- having it embedded turns AETHER-OS from ``an app I use'' into ``the app I live in.''

\textbf{Technology:} \texttt{portable-pty} (Rust) for PTY management, \texttt{xterm.js} (frontend) for terminal rendering, Tauri IPC for streaming output.

\textbf{What you get:} Run commands without context-switching. AI can see terminal output and help you debug. Projects open with a terminal pre-pointed at their directory.

\subsection{1.2 --- System Monitor}

\textbf{What:} Live dashboard showing CPU usage, RAM consumption, disk space, network activity, battery, and top processes.

\textbf{Why:} On an 8\,GB M2 Mac, resource awareness is survival. You need to know if Ollama is eating all your RAM before your Mac freezes. A built-in monitor means you never need Activity Monitor again.

\textbf{Technology:} \texttt{sysinfo} crate (Rust) for hardware metrics, streaming via Tauri events, React charts.

\textbf{What you get:} Real-time visibility into your machine's health. Set alerts for high memory usage. Know exactly when it's safe to run a heavy AI query.

\subsection{1.3 --- Clipboard Manager}

\textbf{What:} System-wide clipboard history. Every text snippet, code block, URL, and image you copy is captured and searchable. Pin important clips, search by keyword, paste with one click.

\textbf{Why:} Developers copy-paste hundreds of times per day. Losing a snippet means re-searching, re-reading, re-copying. A clipboard history turns throwaway copies into a searchable archive.

\textbf{Technology:} Rust backend polling \texttt{NSPasteboard} (macOS), SQLite for history storage, React frontend with search/filter.

\textbf{What you get:} Never lose a copied snippet again. Search paste history by content. Pin frequently used code blocks. Reduces friction between browser, editor, and terminal.

\subsection{1.4 --- Quick Launcher}

\textbf{What:} Spotlight-style launcher (Cmd+Space) that searches across everything: notes, projects, files, commands, bookmarks, AI conversations. One keystroke to jump anywhere or open anything.

\textbf{Why:} The Command Bar (Cmd+K) already exists for navigation. A Quick Launcher extends this to the entire system --- open any file, launch any app, run any command, search any note, all from one input field.

\textbf{Technology:} Extend existing CommandBar component, add file system indexing (Rust \texttt{walkdir}), fuzzy search (\texttt{fuzzy-matcher} crate), app launching via \texttt{open -a}.

\textbf{What you get:} One input to rule them all. No more Finder, no more Spotlight, no more app switching. Type what you want, hit enter, it's open.

% =========================================================
\section{Phase 2: The Knowledge Engine}
% =========================================================

\subsection{2.1 --- Note Editor with Live Preview}

\textbf{What:} Full Markdown editing with live side-by-side preview. Syntax highlighting, wikilink autocomplete, tag autocomplete, image embedding, Mermaid diagram editing.

\textbf{Why:} AETHER-OS currently shows notes in read-only preview. A real knowledge OS lets you write, not just read. Editing is the core loop --- capture $\rightarrow$ edit $\rightarrow$ link $\rightarrow$ recall.

\textbf{Technology:} CodeMirror 6 (lightweight, extensible) or Monaco Editor, custom Markdown parser extensions for wikilinks, live preview via existing MarkdownRenderer.

\textbf{What you get:} Write and edit notes without leaving AETHER-OS. Autocomplete for \texttt{[[wikilinks]]} and \texttt{\#tags}. See your formatting live as you type. Turn AETHER-OS from a reader into a workspace.

\subsection{2.2 --- Backlinks \& Bidirectional Links}

\textbf{What:} When you write \texttt{[[Ekins Work]]} in any note, AETHER-OS automatically tracks the reverse link. A backlinks panel shows every note that references the current note.

\textbf{Why:} Backlinks are the killer feature of Obsidian and Roam. They reveal unexpected connections --- ``oh, I referenced this project in 3 different meeting notes.'' Without backlinks, notes are isolated islands. With them, they're a network.

\textbf{Technology:} Rust backend scanning all vault notes for \texttt{[[wikilink]]} patterns, building a reverse index. Frontend panel showing backlinks with context previews.

\textbf{What you get:} Discover connections you didn't know existed. See which notes reference the current one. Navigate your knowledge graph by relationship, not by folder.

\subsection{2.3 --- Daily Notes \& Quick Capture}

\textbf{What:} A daily note is auto-created each day (e.g., \texttt{2026-07-28.md}). A global hotkey (Cmd+Shift+N) opens a quick-capture input --- type a thought, hit enter, it's appended to today's note.

\textbf{Why:} Friction kills note-taking. If you have to open the app, navigate to a folder, create a file, name it, and start typing --- you won't. Quick capture removes all friction. Daily notes give every thought a home.

\textbf{Technology:} Rust backend creating daily note files, global hotkey registration via Tauri, React modal for quick capture input.

\textbf{What you get:} Capture any thought in under 2 seconds. Every day has a chronological log. Search past days instantly. Build a journaling habit without trying.

\subsection{2.4 --- AI Agent Actions (Write, Run, Automate)}

\textbf{What:} The AI doesn't just answer questions --- it takes actions. ``Create a note titled X'', ``Summarize this URL and save it'', ``Run \texttt{cargo test} in project Y and explain the failures'', ``Add a task to my daily note''.

\textbf{Why:} Chat-only AI is a glorified search engine. Action-capable AI is a collaborator. The difference between ``here's what you should do'' and ``I did it for you'' is the difference between a tool and an OS.

\textbf{Technology:} Tool-calling architecture --- AI responses include structured action intents, Rust backend executes them (file writes, shell commands, note creation) with user approval for destructive actions.

\textbf{What you get:} Tell the AI what you want, it does it. Create notes from conversations. Run commands from chat. Summarize web pages into your vault. The AI becomes an agent, not just a chatbot.

\subsection{2.5 --- Web Clipper \& Research Mode}

\textbf{What:} Paste a URL, AETHER-OS fetches the page, extracts the main content (stripping nav/ads/footer), and saves it as a Markdown note in your vault. Optionally ask the AI to summarize it first.

\textbf{Why:} Research means collecting sources. Currently, you'd copy-paste from a browser manually. A web clipper turns AETHER-OS into a research engine --- capture, summarize, and connect web content without leaving your workspace.

\textbf{Technology:} Rust \texttt{reqwest} for fetching, \texttt{scraper} crate for HTML parsing/content extraction, Markdown conversion, AI summarization via existing Ollama integration.

\textbf{What you get:} Build a personal knowledge base from web content. One click to capture an article. AI summarizes it. It's searchable alongside your own notes. Never lose a useful article again.

% =========================================================
\section{Phase 3: The Productivity Layer}
% =========================================================

\subsection{3.1 --- Task \& Project Management}

\textbf{What:} Kanban board view for tasks. Tasks are extracted from note checkboxes (\texttt{- [ ]}) and can be viewed, dragged, and managed across columns (Todo, In Progress, Done). Tasks link back to their source note.

\textbf{Why:} Tasks scattered across notes are invisible. A Kanban view aggregates them into one actionable surface. You see everything you need to do, prioritize by dragging, and click through to the note for context.

\textbf{Technology:} Rust backend scanning vault for \texttt{- [ ]} and \texttt{- [x]} patterns, extracting task text + source note path. React Kanban board with drag-and-drop.

\textbf{What you get:} All your tasks in one place. Drag to prioritize. Click to see the note context. No separate todo app needed.

\subsection{3.2 --- Calendar \& Reminders}

\textbf{What:} Monthly calendar view showing tasks with due dates, daily notes, and reminders. Desktop notifications for upcoming deadlines. ICS import/export for Google Calendar / Apple Calendar sync.

\textbf{Why:} Time is the missing dimension in note apps. Notes are spatial (they live in folders), but work is temporal (it has deadlines). A calendar bridges this --- see what's due, when, and what your day looks like.

\textbf{Technology:} Rust backend parsing date metadata from note frontmatter, ICS parsing/generation, macOS notification API.

\textbf{What you get:} See your week at a glance. Get notified before deadlines. Sync with your existing calendar. Never miss a deadline because it was buried in a note.

\subsection{3.3 --- Pomodoro \& Focus Mode}

\textbf{What:} Built-in Pomodoro timer (25/5 or custom). Focus mode dims everything except the current note/editor. Timer state persists across sessions. Daily focus time tracked and shown on dashboard.

\textbf{Why:} Productivity isn't just about having tools --- it's about using them consistently. A Pomodoro timer builds focus habits. Focus mode eliminates distractions. Tracked time shows you patterns in your work.

\textbf{Technology:} React timer component, Tauri window management for focus mode, SQLite or JSON for time tracking persistence.

\textbf{What you get:} Build deep work habits. Track how much focused time you actually get. See patterns --- ``I focus best at 10am'' --- and optimize your schedule.

\subsection{3.4 --- Bookmarks \& Pinned Items}

\textbf{What:} Pin frequently used notes, projects, commands, and AI conversations to a sidebar. Quick access without searching. Organize pins into groups.

\textbf{Why:} You have 10--20 things you access daily. Searching for them every time is friction. Pinned items are one click away --- your daily note, your main project, your most-used AI conversation.

\textbf{Technology:} Zustand store with localStorage persistence, React sidebar component, drag-to-reorder.

\textbf{What you get:} Your most-used items are always one click away. No searching, no navigating. Customize your workflow surface.

% =========================================================
\section{Phase 4: The Intelligence Layer}
% =========================================================

\subsection{4.1 --- Auto-Git Versioning for Notes}

\textbf{What:} Every note edit is automatically committed to a local git repository. Browse version history, diff changes, restore previous versions. No manual commits needed.

\textbf{Why:} Notes change over time. Without version history, you lose the evolution of your thinking. Auto-git means every edit is recoverable --- you can always go back. It's like Time Machine for your brain.

\textbf{Technology:} Rust backend running \texttt{git add \&\& git commit} on file change events (via \texttt{notify} crate), git log/diff parsing for history UI, React timeline component.

\textbf{What you get:} Never lose a draft. See how your notes evolved. Restore anything. Your knowledge base has a history, not just a present.

\subsection{4.2 --- AI-Powered Note Suggestions}

\textbf{What:} As you write, the AI suggests related notes you might want to link to. After writing a note, the AI offers to create backlinks, suggest tags, and summarize the note for the dashboard.

\textbf{Why:} Manual linking is tedious. AI suggestions surface connections you'd miss --- ``this note about React hooks is related to your note about state management from last month.'' The system actively helps you build your knowledge graph.

\textbf{Technology:} Semantic similarity via existing vector embeddings, real-time suggestions during editing, AI summarization via Ollama.

\textbf{What you get:} Your knowledge graph grows automatically. Connections you'd never make manually are surfaced. The system gets smarter the more you use it.

\subsection{4.3 --- Conversation Auto-Compaction}

\textbf{What:} Long AI conversations are automatically compressed into summaries when they exceed a token threshold. The summary preserves key facts and decisions. Full transcript is archived but not loaded into context.

\textbf{Why:} Long conversations bloat context windows, slow down AI responses, and consume RAM. Auto-compaction keeps conversations fast and affordable while preserving memory. You can have month-long project conversations without performance degradation.

\textbf{Technology:} Rust backend tracking conversation token count, Ollama summarization when threshold exceeded, summary stored in MemoryStore.

\textbf{What you get:} Month-long AI conversations that stay fast. No more manually starting ``new chat'' because the old one got slow. The AI remembers the summary, you keep the full history.

\subsection{4.4 --- Cross-Module Universal Search}

\textbf{What:} One search bar that searches across everything: notes, projects, AI conversations, memory facts, clipboard history, bookmarks, file system. Results are categorized and ranked by relevance.

\textbf{Why:} Currently, search is fragmented --- semantic search only covers notes, project search only covers projects. Universal search means you type once and find anything, regardless of where it lives.

\textbf{Technology:} Unified search index combining vector embeddings (notes), BM25 keyword search (all text), file system index, and memory store queries. React results UI with category tabs.

\textbf{What you get:} Find anything from one input. No more ``which search do I use?'' One box, everything searchable, results ranked by relevance.

% =========================================================
\section{Phase 5: The Extension Layer}
% =========================================================

\subsection{5.1 --- Plugin System}

\textbf{What:} A plugin architecture allowing custom tools, views, and AI tools to be added without modifying core code. Plugins are JavaScript/TypeScript modules loaded at runtime.

\textbf{Why:} No app anticipates every use case. A plugin system means AETHER-OS grows with your needs --- add a flashcard plugin, a recipe organizer, a workout tracker, whatever you want. The community can build and share plugins.

\textbf{Technology:} Plugin manifest format, sandboxed JS execution (via \texttt{deno\_core} or Web Workers), plugin API exposing vault access, AI access, and UI registration.

\textbf{What you get:} Infinite extensibility without forking. Build custom tools for your workflow. Share them with others. AETHER-OS becomes a platform, not just an app.

\subsection{5.2 --- Export \& Publishing}

\textbf{What:} Export notes to PDF, HTML, or standalone Markdown bundles. Publish a note or folder as a static website (like Obsidian Publish, but self-hosted and free).

\textbf{Why:} Knowledge is meant to be shared. Currently, notes are locked in the vault. Export lets you share with non-AETHER-OS users. Publishing lets you share your knowledge with the world.

\textbf{Technology:} Rust backend for PDF generation, static site generation from Markdown, optional GitHub Pages deployment.

\textbf{What you get:} Share notes as polished PDFs. Publish a knowledge site from your vault. Export everything as a backup. Your knowledge isn't locked in.

\subsection{5.3 --- Sync \& Multi-Device}

\textbf{What:} End-to-end encrypted sync between devices via a simple relay (or direct LAN sync). No cloud account required --- sync key is derived from a password you memorize.

\textbf{Why:} You work on multiple devices --- Mac, maybe a phone eventually. Your knowledge base should be everywhere you are, without trusting a cloud provider with your data.

\textbf{Technology:} Rust sync engine with AES-256-GCM encryption, optional relay server (self-hosted or free tier), conflict resolution via CRDT or last-write-wins.

\textbf{What you get:} Your knowledge base on every device. Encrypted, private, no subscription. Sync happens in the background.

% =========================================================
\section{Priority Matrix}
% =========================================================

\begin{table}[h]
\centering
\small
\begin{tabularx}{\textwidth}{Xccc}
\toprule
\textbf{Feature} & \textbf{Impact} & \textbf{Effort} & \textbf{Priority} \\
\midrule
1.1 Built-in Terminal & Critical & Medium & P0 \\
1.2 System Monitor & High & Low & P0 \\
1.3 Clipboard Manager & High & Medium & P1 \\
1.4 Quick Launcher & High & Low & P1 \\
2.1 Note Editor & Critical & High & P0 \\
2.2 Backlinks & High & Low & P1 \\
2.3 Daily Notes \& Quick Capture & High & Low & P0 \\
2.4 AI Agent Actions & Critical & High & P1 \\
2.5 Web Clipper & Medium & Medium & P2 \\
3.1 Task Management (Kanban) & High & Medium & P1 \\
3.2 Calendar \& Reminders & Medium & High & P2 \\
3.3 Pomodoro \& Focus Mode & Medium & Low & P2 \\
3.4 Bookmarks \& Pinned Items & Medium & Low & P1 \\
4.1 Auto-Git Versioning & High & Medium & P1 \\
4.2 AI Note Suggestions & Medium & Medium & P2 \\
4.3 Conversation Auto-Compaction & High & Low & P1 \\
4.4 Universal Search & High & Medium & P1 \\
5.1 Plugin System & High & Very High & P3 \\
5.2 Export \& Publishing & Medium & Medium & P3 \\
5.3 Sync \& Multi-Device & High & Very High & P3 \\
\bottomrule
\end{tabularx}
\end{table}

% =========================================================
\section{Architecture Impact}
% =========================================================

\textbf{Current Architecture:}
\begin{verbatim}
Frontend (React) <--IPC--> Backend (Rust) --> Ollama + Vault
\end{verbatim}

\textbf{Target Architecture:}
\begin{verbatim}
Frontend (React) <--IPC--> Kernel (Rust)
  |-- Vault Engine (notes, backlinks, daily notes)
  |-- AI Engine (chat, actions, suggestions, compaction)
  |-- Project Engine (scan, git, terminal)
  |-- System Engine (monitor, clipboard, notifications)
  |-- Search Engine (universal: vector + keyword + file)
  |-- Sync Engine (encrypted, multi-device)
  `-- Plugin Runtime (sandboxed JS/TS)
\end{verbatim}

% =========================================================
\section{The North Star}
% =========================================================

AETHER-OS should be the first app you open in the morning and the last you close at night. It should replace:

\begin{itemize}[noitemsep]
  \item \textbf{Activity Monitor} $\rightarrow$ System Monitor
  \item \textbf{Notes app} $\rightarrow$ Vault with full editor
  \item \textbf{Todo app} $\rightarrow$ Task management from notes
  \item \textbf{Calendar app} $\rightarrow$ Built-in calendar with reminders
  \item \textbf{Spotlight} $\rightarrow$ Quick Launcher
  \item \textbf{Clipboard manager} $\rightarrow$ Built-in clipboard history
  \item \textbf{Terminal} $\rightarrow$ Built-in PTY terminal
  \item \textbf{ChatGPT/Claude} $\rightarrow$ Local AI with your context
  \item \textbf{Obsidian} $\rightarrow$ Full knowledge management with AI
\end{itemize}

\vspace{1cm}
\begin{center}
\Large \textit{One app. Your data. Your machine. Your OS.}
\end{center}

\end{document}
